\section{モデル比較}
\subsection{授業内容}
\subsubsection{科目の中でのこのコマの位置づけ}
最終回となるこの回では，これまで後期の授業で扱ってきたさまざまなモデルについて総括し，ベイジアンモデリングの心理学的位置付けについて解説する。
加えて最後の話題提供として，モデル比較について言及する。帰無仮説検定の代わりとして考えるのであれば，区間推定を使った比較が必要であるし，ベイジアンモデリングの観点からはモデル比較になる。ベイズファクターによるモデル比較とそれを実行するためのブリッジサンプリング法，またWAICなど予測的観点から評価する方法があることなどに言及する。
\subsubsection{コマ主題細目}
\begin{description}
	\item[ベイジアンモデリング] 一般的に「モデリング」という観点から，これまでの授業内容だけでなく心理学における研究法としてのその意味や意義を考える。帰無仮説検定のAlternativeとして利用するだけでなく，心理学的メカニズムをより具体的に，緻密に記載するために数学的方法を用いることで，心のメカニズムの理解を深めることができるかもしれない。そこに含まれる仮定や前提について，自覚的に記述する必要があることがモデリングの利点であり，MCMCをはじめとするベイズ統計の技術は，それを可能にしてくれる方法論的補助にすぎない。
	\item[帰無仮説検定の代案] 帰無仮説検定の代わりにモデリングアプローチを取ることの利点はさまざま挙げられるが，p値のように「ここだけ見ておけば良い」というような機械的判断ができる基準がない。むしろそうした機械的判断の弊害が指摘されてきているのであるが，代案としてはどのような基準があるのかはドメイン知識に基づく必要がある。むしろパラメータだけでなくデータの観点から考察できるようになったことや，実質的な値に基づいて考えられることを利点と捉えつつ，判断基準としてのROPEなどについて一瞥する。
		\begin{flushright}
			パラメータ推定かモデル比較かについては，$\to$\citet{Maeda2017}のPp.341--361.
		\end{flushright}
	\item[モデル比較] モデルとその有用性を考えるにあたって，パラメータ推定かモデル比較かという二つのアプローチがあり得る。後者については，階層モデルによるもの，ベイズファクター，WAICが考えられる。ただしWAICについては，渡辺ベイズ理論ともいうべき，より包括的なベイズ理論の枠組みで捉え直す必要があり，心理学研究にこの枠組みがどれほど有用かについては，いまだに結論が出ていないところでもある。ここでは概略的にその特徴に触れるにとどまり，受講生諸君の今後の活躍に期待したい。
		\begin{flushright}
			WAICについては$\to$\citet{浜田201912}が丁寧である。
		\end{flushright}
\end{description}
\subsubsection{キーワード}
\begin{itemize}
	\item ROPE
	\item ベイズファクター
	\item サヴェージ・ディッキー法
	\item ブリッジサンプリング
\end{itemize}
\subsection{授業情報}
\paragraph{コマの展開方法}
講義/遠隔可/演習
